\documentclass[12 pt]{article}
\usepackage[spanish]{babel}
\usepackage{amssymb}
%opening
\title{Trabajo Practico 1}
\author{Deppen Nahuel}
\date{}

%------------
\usepackage{setspace}
\doublespacing

\begin{document}

\maketitle



\section{${v}_{j}\in S.E({v}_{1},...,{v}_{n}), j=1,...,n$}
\underline{\textbf{Demos:}}  \\
\\
  Veamos que por definición tenemos que $v={a}_{1}{v}_{1}+{a}_{2}{v}_{2}+...+{a}_{j}{v}_{j}+...+{a}_{n}{v}_{n}$
 Tomando ${a}_{j}=1$ y ${a}_{i}=0, i=1,..,j-1,j+1,...n$ \\ Tenemos:\\
 
  \begin{center} $v=0{v}_{1}+0{v}_{2}+...+{v}_{j}+....+0{v}_{n}={v}_{j}$
\end{center}

 Luego ${v}_{j} $es una combinacion lineal de ${v}_{1},...,{v}_{n}\in V, j=1,...,n$
 
 \begin{center}
 	$ \therefore {v}_{j} \in  S.E({v}_{1},...,{v}_{n}), j=1,...,n$
 \end{center} 
 
 

 \newpage
\section{$S.E({v}_{1},...,{v}_{n})$ es un subespacio de V}
\underline {\textbf{Demos:}}\\
1) Veamos que podemos escribir a 0 (cero), como una CL, con

${v}_{j}\neq0, {a}_{j}=0, j=1,...,n$:

\begin{center}
	${a}_{1}{v}_{1}+{a}_{2}{v}_{2}+...+{a}_{n}{v}_{n}=
	0{v}_{1}+0{v}_{2}+....+0{v}_{n}={v}_{j}=0 $ 
	
\end{center}

 
Si existe algun ${v}_{j}=0$ entonces tomando $ {a}_{i}=0, i=1,...,j-1,j+1,...n$

\begin{center}
	$v=0{v}_{1}+...+{v}_{j}+...+{v}_{n}={v}_{j}=0$
\end{center}
\begin{center}
$ \therefore 0 \in S.E({v}_{1},...,{v}_{n})$
\end{center}




2)Dados $ v \in S.E({v}_{1},...,{v}_{n})$ y
$w \in S.E({v}_{1},...,{v}_{n})$


Dados ${a}_{i} \in \Bbbk $ y ${b}_{i} \in \Bbbk , i=1,...,n$

\begin{center}
	
$v={a}_{1}{v}_{1}+{a}_{2}{v}_{2}+...+{a}_{n}{v}_{n} \in  S.E({v}_{1},...,{v}_{n})$

$w={b}_{1}{v}_{1}+{b}_{2}{v}_{2}+...+{b}_{n}{v}_{n} \in  S.E({v}_{1},...,{v}_{n})$
 


\end{center}



Tenemos:

\begin{center}
	$v+w=({a}_{1}{v}_{1}+{a}_{2}{v}_{2}+...+{a}_{n}{v}_{n})+({b}_{1}{v}_{1}+{b}_{2}{v}_{2}+...+{b}_{n}{v}_{n})=
	{c}_{1}{v}_{1}+{c}_{2}{v}_{2}+...+{c}_{n}{v}_{n} $

	
	con ${c}_{i}={b}_{i}+{a}_{i}, i=1...n$ y ${c}_{i} \in \Bbbk$
\end{center}
 
 $ \therefore v+w \in S.E({v}_{1},...,{v}_{n})$
 
 \newpage
 
 3) Dado $v \in S.E({v}_{1},...,{v}_{n})$ y $\alpha \in \Bbbk $:
 
 \begin{center}
 	
 	$ \alpha v= \alpha ({a}_{1}{v}_{1}+{a}_{2}{v}_{2}+...+{a}_{n}{v}_{n})=
 	{\alpha a}_{1}{v}_{1}+{\alpha a}_{2}{v}_{2}+...+{\alpha a}_{n}{v}_{n}$
 	
 	Con $ {\alpha a}_{i} \in \Bbbk , i=1,...,n$
\end{center}

$ \therefore \alpha v \in S.E({v}_{1},...,{v}_{n})$


\begin{center} 
	
	$ \therefore $ Por 1), 2), y 3) $ S.E({v}_{1},...,{v}_{n}) $ es un subespacio de V 
\end{center}



  \section{Si $U \subset V$ es un subespacio tal que ${v}_{1},...,{v}_{n} \in U$, luego $S.E({v}_{1},...,{v}_{n}) \subset U$ }

\underline{\textbf{Demos:}}

Como U es un s.e de V, sabemos que $0 \in U$ y , la suma y  multiplicación por escalar, son cerradas para todo ${v}_{1},...,{v}_{n} \in U. $

Ahora, ${v}_{1},...,{v}_{n} \in U$  solo si ${v}_{1},...,{v}_{n} \in V$, como vimos en el punto anterior dados ${v}_{1},...,{v}_{n} \in V$ sabemos que  $ S.E({v}_{1},...,{v}_{n}) $ es un subespacio de V. 

Luego, como ${v}_{1},...,{v}_{n} \in U $ y $ S.E({v}_{1},...,{v}_{n}) \subset V $ entonces   $ S.E({v}_{1},...,{v}_{n}) \subset U $

\end{document}

